\documentclass[a4paper]{article}

\usepackage{graphicx}
\usepackage{amsmath,amssymb}
\usepackage{minted}
\usepackage{multirow}
\usepackage{float}
\usepackage{algorithm}
\usepackage{algpseudocode}
\usepackage{todonotes}
\usepackage{forest}
\usepackage{xstring}
\usepackage{xcolor}
\usepackage[most]{tcolorbox}
\usepackage{xparse}
\usepackage{natbib}

\usetikzlibrary{graphs,graphdrawing}
\usegdlibrary{layered}

% for external graphviz
\input{/home/donkarlo/Dropbox/repo/nd_language_ai_project/src/nd_language_ai/formal/computer/programming/tex/latex/code_clips/external_graphviz.tex}

% for tags
\input{/home/donkarlo/Dropbox/repo/nd_language_ai_project/src/nd_language_ai/formal/computer/programming/tex/latex/parts/control_sequence/kind/semtag/semtag.tex}

% for links
\input{/home/donkarlo/Dropbox/repo/nd_language_ai_project/src/nd_language_ai/formal/computer/programming/tex/latex/code_clips/seehere.tex}

\title{}
\date{}
\author{Mohammad Rahmani}

\begin{document}

   \maketitle

   \section{Text 1}

      Früher hat Frau Berger oft allein zu Hause gegessen – ohne Appetit. Seit sie beim wöchentlichen Mittagstisch im Nachbarschaftshaus mitmacht, ist das anders. Sie sagt: „Ich koche wieder gern, wenn ich weiß, dass ich nicht allein esse.“ Hier treffen sich Menschen aus dem Viertel, kochen gemeinsam und tauschen sich aus. Das Projekt zeigt: Zusammen zu essen macht nicht nur satt, sondern auch glücklich.

   \section{Text 2}

      Der 66-jährige Herr Albrecht leidet seit Jahren unter Rückenschmerzen. Medikamente halfen nur wenig. Durch Zufall nahm er an einem Bewegungsprogramm seiner Krankenkasse teil – mit erstaunlichem Erfolg. Tägliche Spaziergänge, leichte Übungen und gezielte Atemtechniken verbesserten seinen Alltag enorm. „Ich fühle mich aktiver und brauche weniger Tabletten“, sagt er. Für ihn ist klar: Bewegung ist die beste Medizin.

   \section{Text 3}

      Smartphones, Tablets, Laptops – in vielen Haushalten laufen mehrere Geräte gleichzeitig. Familie Müller hat festgestellt, dass alle Familienmitglieder immer unkonzentrierter wurden. Deshalb beschlossen sie, abends eine „Technikpause“ einzulegen: kein WLAN, keine Bildschirme, kein Multitasking. Anfangs war das ungewohnt, aber inzwischen freuen sich alle auf diese ruhige Zeit. Sie spielen Brettspiele, lesen oder reden einfach miteinander – ohne ständige Ablenkung.

   \section{Text 4}

      In einem alten Gebäude in der Innenstadt hat sich ein neues Freizeitkonzept etabliert: kleine, flexible Räume, die stundenweise gemietet werden können. Hier treffen sich Leute zum Basteln, Musizieren oder einfach zum Reden. Die Idee entstand, weil viele in kleinen Wohnungen leben und keinen Platz für Hobbys oder Besuch haben. Die Nutzer schätzen die Abwechslung und die Möglichkeit, neue Menschen kennenzulernen.

   \section{Text 5}

      Kochen war für Jana nie ein Thema – Fertiggerichte waren ihr Alltag. Nach einem Burnout riet ihr ihre Therapeutin zu einer Kochgruppe für Stressbetroffene. Dort lernte sie einfache Rezepte kennen, aber auch, wie das Kochen selbst beruhigen kann. „Ich bin konzentriert, rieche die Zutaten, habe ein Ziel – das tut mir gut“, sagt sie. Das gemeinsame Essen ist dabei nur ein Bonus.

      \bibliographystyle{plainnat}
      \bibliography{/home/donkarlo/Dropbox/repo/nd_shared_research/bibliography/bibliography.bib}

\end{document}
